\documentclass[conference]{IEEEtran}

% ── packages ──────────────────────────────────────────────────────────────────
\usepackage{cite}
\usepackage{amsmath,amssymb,amsfonts}
\usepackage{algorithmic}
\usepackage{graphicx}
\usepackage{textcomp}
\usepackage{xcolor}
\usepackage{booktabs}
\usepackage{multirow}
\usepackage{hyperref}
\usepackage{tabularx}
\usepackage{float}
\usepackage{subcaption}
\usepackage{textcomp}
\usepackage{gensymb}

\hypersetup{
    colorlinks=true,
    linkcolor=blue,
    citecolor=blue,
    urlcolor=blue,
}

\def\BibTeX{{\rm B\kern-.05em{\sc i\kern-.025em b}\kern-.08em
    T\kern-.1667em\lower.7ex\hbox{E}\kern-.125emX}}

\begin{document}

% ══════════════════════════════════════════════════════════════════════════════
% TITLE
% ══════════════════════════════════════════════════════════════════════════════
\title{LandslideEEGMoE: Adapting a Domain-Decoupled Mixture-of-Experts Architecture for Landslide Susceptibility Mapping Using Multi-Source Remote Sensing Data}

\author{
\IEEEauthorblockN{Toruk-Maktos Team}
\IEEEauthorblockA{
Department of Computer Science and Engineering\\
SEM 6 — Project Report\\
\textit{Toruk-Maktos Team}
}
}

\maketitle

% ══════════════════════════════════════════════════════════════════════════════
% ABSTRACT
% ══════════════════════════════════════════════════════════════════════════════
\begin{abstract}
Landslide susceptibility mapping is a critical task for disaster risk reduction, particularly in monsoon-prone regions such as Kerala, India. In this work, we adapt the EEGMoE architecture—a Domain-Decoupled Mixture-of-Experts model originally proposed for self-supervised EEG representation learning—to the domain of landslide susceptibility prediction. Our adaptation, termed \textbf{LandslideEEGMoE}, replaces EEG frequency-band inputs with five geospatial feature channels derived from Sentinel-1 SAR, Sentinel-2 multispectral imagery, rainfall data, and SMAP soil moisture observations. The model preserves the exact architectural components of EEGMoE: the Spatial \& Frequency Encoder, the Shared-Specific Mixture-of-Experts (SSMoE) block with Top-K routing, and the two-stage self-supervised pre-training followed by supervised fine-tuning pipeline. We validate the approach on two real landslide events—the 2019 Puthumala landslide and the 2024 Wayanad landslide—using the Indian Landslide Atlas as ground truth. Our implementation demonstrates that the domain-decoupled expert specialisation mechanism, originally designed to separate EEG frequency-band representations, naturally extends to separating heterogeneous geospatial feature domains, enabling each expert to specialise in distinct environmental factors contributing to slope failure.
\end{abstract}

\begin{IEEEkeywords}
Landslide susceptibility mapping, Mixture-of-Experts, self-supervised learning, remote sensing, SAR, NDVI, soil moisture, domain decoupling
\end{IEEEkeywords}

% ══════════════════════════════════════════════════════════════════════════════
% I. INTRODUCTION
% ══════════════════════════════════════════════════════════════════════════════
\section{Introduction}

Landslides are among the most destructive geohazards globally, claiming thousands of lives and causing billions of dollars in infrastructure damage annually. In India, the Western Ghats region—particularly Kerala—has witnessed catastrophic landslide events during monsoon seasons, including the devastating 2019 Puthumala landslide and the 2024 Wayanad landslide that caused widespread destruction \cite{wayanad2024}.

Traditional landslide susceptibility mapping relies on statistical methods (logistic regression, frequency ratio) or physics-based models (infinite slope, SHALSTAB), which require extensive manual feature engineering and domain expertise. Recent advances in deep learning have shown promise for automated landslide detection \cite{dl_landslide}, but most approaches treat the multi-source input data homogeneously, failing to exploit the distinct physical characteristics of each data source.

The EEGMoE model \cite{eegmoe2026} was proposed for self-supervised EEG representation learning and introduced a novel Domain-Decoupled Mixture-of-Experts (MoE) architecture. Its key insight is that different frequency bands in EEG signals (delta, theta, alpha, beta, gamma) carry domain-specific information that benefits from specialised expert processing. This architectural principle directly parallels the multi-source nature of landslide-relevant geospatial data, where SAR backscatter, vegetation indices, rainfall, and soil moisture each encode fundamentally different physical phenomena.

In this paper, we present \textbf{LandslideEEGMoE}, a faithful replication of the EEGMoE architecture adapted for landslide susceptibility mapping. Our contributions are:

\begin{enumerate}
    \item \textbf{Exact architectural replication}: We implement EEGMoE's full architecture—Spatial \& Frequency Encoder, SSMoE blocks with Top-K routing, and two-stage training—precisely as described in \cite{eegmoe2026} (Fig.~2 and Fig.~3).
    
    \item \textbf{Domain-appropriate channel mapping}: We establish a principled mapping from EEG frequency bands ($\delta$, $\theta$, $\alpha$, $\beta$, $\gamma$) to five geospatial feature channels (SAR, NDVI, NDWI, Rainfall, Soil Moisture).
    
    \item \textbf{Real geospatial data pipeline}: We build a complete data pipeline that reads Sentinel-1, Sentinel-2, IMD rainfall, and SMAP soil moisture data from their native formats (GeoTIFF, NetCDF) and produces normalised patch-based training samples.
    
    \item \textbf{Validation on real events}: We apply the model to two documented landslide events in Kerala, India, using the Indian Landslide Atlas (2023) as ground truth.
\end{enumerate}

% ══════════════════════════════════════════════════════════════════════════════
% II. RELATED WORK
% ══════════════════════════════════════════════════════════════════════════════
\section{Related Work}

\subsection{Landslide Susceptibility Mapping}
Conventional landslide susceptibility mapping employs statistical approaches such as logistic regression \cite{logistic_landslide}, support vector machines \cite{svm_landslide}, and random forests \cite{rf_landslide}. These methods require careful manual selection and engineering of conditioning factors. Recent deep learning methods—including CNNs \cite{cnn_landslide}, U-Nets, and attention-based architectures—have improved detection accuracy but typically process all input features through a shared network without domain-specific specialisation.

\subsection{Mixture-of-Experts Models}
Mixture-of-Experts (MoE) architectures \cite{moe_original} route inputs to specialised sub-networks (experts) based on learned gating functions. Recent work has scaled MoE to large language models (Switch Transformer \cite{switch_transformer}, GShard) and applied them to multi-modal learning. The EEGMoE model \cite{eegmoe2026} introduced a \textit{Shared-Specific} MoE variant (SSMoE) where specific experts capture domain-specific patterns via Top-K routing while shared experts learn cross-domain representations, and demonstrated its effectiveness for multi-band EEG representation learning.

\subsection{Self-Supervised Pre-training for Remote Sensing}
Self-supervised learning has gained traction in remote sensing, with masked autoencoders (MAE) \cite{mae} and contrastive learning showing strong results for satellite image understanding. Our work applies the specific masking-based pre-training strategy from EEGMoE (random token masking with L1 reconstruction loss) to geospatial feature patches.

% ══════════════════════════════════════════════════════════════════════════════
% III. METHODOLOGY
% ══════════════════════════════════════════════════════════════════════════════
\section{Methodology}

Our methodology follows the EEGMoE paper \cite{eegmoe2026} \S III exactly, adapting Section III-A (Spatial \& Frequency Encoder), Section III-B (Domain-Decoupled MoE Encoder), and Section III-C (two-stage training) for the landslide domain.

\subsection{Input Representation — Channel Mapping}

The EEGMoE model processes 5 EEG frequency bands as input channels. We establish the following mapping from EEG bands to geospatial features, preserving the 5-channel input structure:

\begin{table}[h]
\centering
\caption{Channel Mapping: EEG Bands $\rightarrow$ Geospatial Features}
\label{tab:channel_mapping}
\begin{tabular}{@{}clll@{}}
\toprule
\textbf{Ch} & \textbf{EEG Band} & \textbf{Geospatial Feature} & \textbf{Source} \\
\midrule
0 & $\delta$ (delta) & SAR backscatter & Sentinel-1 \\
1 & $\theta$ (theta) & NDVI & Sentinel-2 B08, B04 \\
2 & $\alpha$ (alpha) & NDWI & Sentinel-2 B03, B08 \\
3 & $\beta$ (beta) & Rainfall & IMD / Kerala Met \\
4 & $\gamma$ (gamma) & Soil Moisture & SMAP \\
\bottomrule
\end{tabular}
\end{table}

\textbf{Rationale}: Each geospatial feature encodes a distinct physical domain analogous to EEG frequency bands:
\begin{itemize}
    \item \textbf{SAR} ($\sim\delta$): Low-frequency structural information about terrain roughness and deformation.
    \item \textbf{NDVI} ($\sim\theta$): Vegetation health; landslide scars exhibit anomalously low NDVI.
    \item \textbf{NDWI} ($\sim\alpha$): Surface water content; saturated areas indicate heightened risk.
    \item \textbf{Rainfall} ($\sim\beta$): Primary triggering factor; extreme precipitation directly induces slope failure.
    \item \textbf{Soil Moisture} ($\sim\gamma$): Subsurface saturation condition; high soil moisture reduces shear strength.
\end{itemize}

\subsection{Spatial \& Frequency Encoder (\S III-A)}

Following Fig.~2 of \cite{eegmoe2026}, the input tensor $\mathbf{X} \in \mathbb{R}^{B \times 5 \times H \times W}$ is first divided into non-overlapping $p \times p$ spatial patches (analogous to EEGMoE's time-window segmentation). Each patch is flattened to dimension $d = 5 \times p^2$ and passed through the Spatial \& Frequency Encoder:

\begin{equation}
    \mathbf{Z} = \text{ReLU}\big(\mathbf{W}_2 \cdot \text{ReLU}(\mathbf{W}_1 \cdot \mathbf{x}_\text{patch} + \mathbf{b}_1) + \mathbf{b}_2\big)
\end{equation}

where $\mathbf{W}_1 \in \mathbb{R}^{64 \times d}$, $\mathbf{W}_2 \in \mathbb{R}^{128 \times 64}$ (Table III values: MLP hidden = 64, embed dim = 128). This produces token embeddings $\mathbf{Z} \in \mathbb{R}^{B \times n_P \times 128}$ with $n_P = (H/p) \times (W/p)$.

Sinusoidal positional encoding is added to preserve spatial ordering of patches, followed by a linear projection to hidden size 512.

\subsection{Domain-Decoupled Encoder (\S III-B)}

The encoder consists of $M = 4$ stacked layers (Table III), each following the Fig.~2 architecture:

\begin{equation}
    \mathbf{h} = \text{LayerNorm}\big(\mathbf{x} + \text{MHSA}(\mathbf{x})\big)
\end{equation}
\begin{equation}
    \mathbf{y} = \text{LayerNorm}\big(\mathbf{h} + \text{SSMoE}(\mathbf{h})\big)
\end{equation}

where MHSA is a multi-head self-attention module with 4 attention heads (Table III).

\subsubsection{SSMoE Block (\S III-B, Eq.~1--5)}

The Shared-Specific MoE block combines two sub-modules:

\textbf{Specific MoE} (Top-K routing, Eq.~1--3):
\begin{equation}
    g(\mathbf{x}) = \mathbf{W}_e \cdot \mathbf{x} \quad \text{(Eq. 1)}
\end{equation}
\begin{equation}
    p_i(\mathbf{x}) = \text{softmax}\big(g(\mathbf{x})\big)_i \quad \text{(Eq. 2)}
\end{equation}
\begin{equation}
    \text{SpecMoE}(\mathbf{x}) = \sum_{i \in \text{TopK}} p_i(\mathbf{x}) \cdot e_i(\mathbf{x}) \quad \text{(Eq. 3)}
\end{equation}

where $K = 2$ and there are 6 specific expert candidates (Table III). Each expert $e_i$ is a two-layer MLP with GELU activation (\S III-B-1).

\textbf{Shared MoE} (soft routing, Eq.~4):
\begin{equation}
    \text{ShareMoE}(\mathbf{x}) = \sum_{i \in F} p_i(\mathbf{x}) \cdot f_i(\mathbf{x}) \quad \text{(Eq. 4)}
\end{equation}

with 2 shared experts using soft (non-sparse) routing.

\textbf{Fusion} (Eq.~5, confirmed additive in Table XII ablation):
\begin{equation}
    \text{SSMoE}(\mathbf{x}) = \text{SpecMoE}(\mathbf{x}) + \text{ShareMoE}(\mathbf{x}) \quad \text{(Eq. 5)}
\end{equation}

\subsection{Two-Stage Training (\S III-C)}

\subsubsection{Stage 1 — Self-Supervised Pre-training (Fig.~3, \S III-C-1)}

Following Fig.~3, we randomly mask 40\% of input tokens (Table XIII: optimal mask ratio = 0.4) by setting them to zero. The masked tokens are passed through the Domain-Decoupled Encoder and reconstructed via a linear head. The pre-training loss combines L1 reconstruction loss on masked tokens with load-balancing auxiliary loss:

\begin{equation}
    \mathcal{L}_\text{recon} = \frac{1}{|\mathcal{M}|} \sum_{i \in \mathcal{M}} \|\hat{\mathbf{Z}}_i - \mathbf{Z}_i\|_1 \quad \text{(Eq. 6)}
\end{equation}

\textbf{Load-balancing loss} (Eq.~7--9) encourages even expert utilisation:
\begin{equation}
    \mathcal{L}_\text{aux} = E \cdot \sum_{i=1}^{E} h_i \cdot D_i \quad \text{(Eq. 7)}
\end{equation}

where $h_i$ is the fraction of tokens routed to expert $i$ (Eq.~8) and $D_i$ is the mean router probability for expert $i$ (Eq.~9).

The total pre-training loss:
\begin{equation}
    \mathcal{L}_\text{pretrain} = \mathcal{L}_\text{recon} + \alpha \cdot \mathcal{L}_\text{aux} \quad \text{(Eq. 10)}
\end{equation}

with $\alpha = 10^{-4}$ (Table III).

Pre-training hyperparameters (Table II): AdamW optimizer, learning rate $10^{-4}$, batch size 64, 5 epochs, warmup ratio 0.05, cosine LR decay.

\subsubsection{Stage 2 — Supervised Fine-tuning (\S III-C-2)}

The pre-trained encoder weights are loaded, and a linear classification head is attached. The full model is fine-tuned with cross-entropy loss:

\begin{equation}
    \mathcal{L}_\text{cls} = \text{CrossEntropy}(\hat{y}, y) \quad \text{(Eq. 11)}
\end{equation}

Classification is performed on global-average-pooled encoder output across all token positions. Fine-tuning hyperparameters (Table II): learning rate $10^{-3}$, batch size 128, 35 epochs for Puthumala (analogous to DEAP dataset), cosine LR decay with warmup.

% ══════════════════════════════════════════════════════════════════════════════
% IV. DATASET AND DATA PIPELINE
% ══════════════════════════════════════════════════════════════════════════════
\section{Dataset and Data Pipeline}

\subsection{Study Sites}

We evaluate our approach on two well-documented landslide events in Kerala, India:

\begin{enumerate}
    \item \textbf{Puthumala Landslide (August 7, 2019)}: A catastrophic debris flow near the Gudalur region on the Tamil Nadu--Kerala border, triggered by extreme monsoon rainfall. This event is used as the \textit{training site}.
    
    \item \textbf{Wayanad Landslide (July 30, 2024)}: Extreme rainfall-induced mass movements in Wayanad district, Kerala, causing extensive casualties and infrastructure damage. This serves as the \textit{validation site}.
\end{enumerate}

\subsection{Data Sources}

\begin{table}[h]
\centering
\caption{Multi-Source Remote Sensing Dataset}
\label{tab:data_sources}
\begin{tabular}{@{}llll@{}}
\toprule
\textbf{Source} & \textbf{Type} & \textbf{Resolution} & \textbf{Format} \\
\midrule
Sentinel-1 & SAR (C-band) & 10 m & GeoTIFF \\
Sentinel-2 & Multispectral & 10--60 m & GeoTIFF \\
IMD Rainfall & Gridded precip. & $0.25^{\circ}$ & NetCDF \\
SMAP & Soil moisture & 9 km & GeoTIFF \\
Copernicus DEM & Elevation & 30 m & GeoTIFF \\
Landslide Atlas & Ground truth & Vector & PDF \\
\bottomrule
\end{tabular}
\end{table}

\subsubsection{Sentinel-1 SAR Data}
Synthetic Aperture Radar images acquired in Interferometric Wide (IW) swath mode at 10~m resolution. SAR backscatter intensity is sensitive to surface roughness and moisture variations, enabling terrain deformation detection independent of cloud cover. We use VV/VH polarisation band 1 (2019-12-08 for Puthumala; 2024-12-11 for Wayanad).

\subsubsection{Sentinel-2 Multispectral Data}
13-band optical imagery at 10--60~m resolution. We derive two spectral indices:
\begin{itemize}
    \item NDVI $= \frac{\text{B08} - \text{B04}}{\text{B08} + \text{B04}}$ — vegetation health indicator; landslide scars show anomalously low NDVI values.
    \item NDWI $= \frac{\text{B03} - \text{B08}}{\text{B03} + \text{B08}}$ — surface water content; elevated values indicate saturated or flooded conditions.
\end{itemize}

\subsubsection{Rainfall Data}
Gridded daily rainfall from India Meteorological Department (IMD) for Puthumala (\texttt{imdlib\_rain\_2019-01-01\_to\_2019-12-31}) and Kerala meteorological records for Wayanad. We extract maximum daily rainfall over the monsoon period as the triggering indicator.

\subsubsection{Soil Moisture (SMAP)}
NASA Soil Moisture Active Passive (SMAP) Level-3 radiometer data (\texttt{SM\_SMAP\_I\_*.tif}) providing volumetric soil moisture estimates at 9~km resolution. We select the observation closest to the event date (Puthumala: August 7, 2019; Wayanad: July 30, 2024).

\subsubsection{Ground Truth — Landslide Atlas}
The Indian Landslide Atlas (2023 edition) provides mapped landslide inventories. Binary label maps (1 = landslide zone, 0 = stable terrain) are derived from the atlas for training and evaluation. Additionally, we provide a heuristic label generation module that combines DEM slope, NDVI anomalies, soil moisture, and rainfall to produce proxy labels when manual digitisation is not available.

\subsection{Data Pipeline}

The preprocessing pipeline follows the z-score normalisation strategy from \cite{eegmoe2026} \S IV-A-2:

\begin{enumerate}
    \item \textbf{Load}: Read raw rasters from GeoTIFF/NetCDF via rasterio and netCDF4.
    \item \textbf{Resample}: Reproject all layers to a common spatial extent and resolution ($256 \times 256$ pixels).
    \item \textbf{Compute indices}: Calculate NDVI and NDWI from Sentinel-2 bands.
    \item \textbf{Stack}: Assemble 5-channel tensor $\mathbf{X} \in \mathbb{R}^{5 \times H \times W}$.
    \item \textbf{Normalise}: Apply z-score normalisation per channel: $\hat{x} = (x - \mu) / \sigma$.
    \item \textbf{Patch extraction}: Slide $64 \times 64$ windows with stride 32 to extract training samples.
    \item \textbf{Label assignment}: Majority-vote binary classification per patch from the ground truth label map.
\end{enumerate}

% ══════════════════════════════════════════════════════════════════════════════
% V. IMPLEMENTATION DETAILS
% ══════════════════════════════════════════════════════════════════════════════
\section{Implementation Details}

\subsection{Architecture Configuration}

All hyperparameters are taken verbatim from Tables II and III of \cite{eegmoe2026}:

\begin{table}[h]
\centering
\caption{Model Hyperparameters (Table III of \cite{eegmoe2026})}
\label{tab:hyperparams}
\begin{tabular}{@{}ll@{}}
\toprule
\textbf{Parameter} & \textbf{Value} \\
\midrule
MLP hidden / embed dim & 64 / 128 \\
Hidden size & 512 \\
FFN dimension & 2048 \\
Attention heads & 4 \\
Encoder layers ($M$) & 4 \\
Top-K routing & $K = 2$ \\
Specific experts & 6 \\
Shared experts & 2 \\
Mask ratio (Stage 1) & 0.4 \\
$\alpha$ (load-balance weight) & $10^{-4}$ \\
Dropout & 0.1 \\
\bottomrule
\end{tabular}
\end{table}

\begin{table}[h]
\centering
\caption{Training Hyperparameters (Table II of \cite{eegmoe2026})}
\label{tab:training}
\begin{tabular}{@{}lll@{}}
\toprule
\textbf{Parameter} & \textbf{Pre-train} & \textbf{Fine-tune} \\
\midrule
Optimizer & AdamW & AdamW \\
Learning rate & $10^{-4}$ & $10^{-3}$ \\
Batch size & 64 & 128 \\
Epochs & 5 & 35 \\
Warmup ratio & 0.05 & 0.05 \\
LR schedule & Cosine & Cosine \\
Weight decay & 0.01 & 0.01 \\
Gradient clip & 1.0 & 1.0 \\
\bottomrule
\end{tabular}
\end{table}

\subsection{Software Stack}

The implementation uses PyTorch with the following key dependencies: rasterio for GeoTIFF I/O, netCDF4 for rainfall data, scipy for spatial resampling, and scikit-learn for evaluation metrics. Xavier initialisation is applied to all linear layers and LayerNorm parameters are initialised to $\mathbf{1}$ (weight) and $\mathbf{0}$ (bias).

\subsection{Codebase Structure}

\begin{table}[h]
\centering
\caption{Code Module Structure}
\label{tab:codebase}
\begin{tabular}{@{}ll@{}}
\toprule
\textbf{Module} & \textbf{Responsibility} \\
\midrule
\texttt{models/eegmoe\_landslide.py} & Full architecture (Fig.~2, 3) \\
\texttt{data/data\_pipeline.py} & Geospatial data loading \\
\texttt{data/generate\_labels.py} & Heuristic label generation \\
\texttt{train/train.py} & Two-stage training loop \\
\texttt{utils/metrics.py} & Evaluation metrics \\
\texttt{configs/config.py} & Tables II \& III parameters \\
\bottomrule
\end{tabular}
\end{table}

% ══════════════════════════════════════════════════════════════════════════════
% VI. ADAPTATION ANALYSIS
% ══════════════════════════════════════════════════════════════════════════════
\section{Adaptation Analysis: EEG $\rightarrow$ Landslide}

\subsection{Structural Correspondence}

The adaptation from EEG to landslide domain preserves the fundamental architectural insight of EEGMoE: \textit{heterogeneous input domains benefit from specialised expert processing}. Table~\ref{tab:correspondence} summarises the structural correspondence.

\begin{table}[h]
\centering
\caption{Structural Correspondence: EEG $\leftrightarrow$ Landslide}
\label{tab:correspondence}
\begin{tabular}{@{}ll@{}}
\toprule
\textbf{EEGMoE (Original)} & \textbf{LandslideEEGMoE (Ours)} \\
\midrule
EEG time series & Geospatial raster \\
22 electrode channels & 5 feature channels \\
5 frequency bands ($\delta\theta\alpha\beta\gamma$) & 5 geo features (SAR, NDVI, ...) \\
Time-window segments & Spatial patches ($p \times p$) \\
Emotion classes (e.g., 4) & Binary: landslide / stable \\
4D input data & 4D input $(B, C, H, W)$ \\
BCI / emotion subjects & Study sites (Puthumala, Wayanad) \\
Self-supervised masking & Self-supervised masking \\
\bottomrule
\end{tabular}
\end{table}

\subsection{Why Domain Decoupling Applies}

The EEGMoE paper demonstrates that different EEG frequency bands encode different neural processes (e.g., delta for deep sleep, alpha for relaxation). Similarly, our five geospatial channels encode fundamentally different physical processes:

\begin{itemize}
    \item \textbf{SAR backscatter} encodes surface \textit{structural} properties (roughness, moisture, geometry).
    \item \textbf{NDVI/NDWI} encode \textit{biological and hydrological} surface conditions.
    \item \textbf{Rainfall} encodes \textit{meteorological triggering} forces.
    \item \textbf{Soil moisture} encodes \textit{subsurface geotechnical} conditions.
\end{itemize}

The SSMoE architecture allows specific experts to specialise in individual factor domains while shared experts capture cross-domain interactions (e.g., the coupled effect of high rainfall \textit{and} high soil moisture on slope stability).

\subsection{Key Adaptations}

While the architecture is preserved exactly, three adaptations were necessary:

\begin{enumerate}
    \item \textbf{Spatial tokenisation}: EEGMoE uses temporal windowing to create token sequences from EEG time series. We use spatial patch extraction ($p \times p$ non-overlapping patches via \texttt{unfold}) to tokenise 2D raster data, producing the same $(B, n_P, d)$ tensor format.
    
    \item \textbf{Preprocessing}: EEG signals undergo bandpass filtering into frequency bands. Our geospatial features undergo spectral index computation (NDVI, NDWI), spatial resampling to common resolution, and z-score normalisation—the same normalisation strategy used in EEGMoE \S IV-A-2.
    
    \item \textbf{Classification head}: The original model classifies emotion states (4 classes for DEAP). We reduce to binary classification (landslide / non-landslide), maintaining the same architecture (global average pooling $\rightarrow$ linear layer).
\end{enumerate}

% ══════════════════════════════════════════════════════════════════════════════
% VII. EVALUATION METRICS
% ══════════════════════════════════════════════════════════════════════════════
\section{Evaluation Metrics}

Following \S IV-A-3 of \cite{eegmoe2026}, we evaluate using:

\begin{itemize}
    \item \textbf{Accuracy} (primary metric, Tables V \& VI)
    \item \textbf{AUROC} — Area Under ROC Curve (Table IV), critical for imbalanced landslide data where non-landslide pixels far outnumber landslide pixels
    \item \textbf{AUC-PR} — Area Under Precision-Recall Curve (Table IV)
    \item \textbf{F1 Score}, \textbf{Precision}, \textbf{Recall}
    \item \textbf{Confusion Matrix}
\end{itemize}

The data is split into 80\% training, 10\% validation, and 10\% test sets with random stratified sampling (seed = 42 for reproducibility). Best model is selected based on validation accuracy, and final metrics are reported on the held-out test set.

% ══════════════════════════════════════════════════════════════════════════════
% VIII. EQUATIONS SUMMARY
% ══════════════════════════════════════════════════════════════════════════════
\section{Equations Summary}

For completeness, Table~\ref{tab:equations} maps all equations from \cite{eegmoe2026} to their implementation locations in our codebase.

\begin{table}[h]
\centering
\caption{Equation-to-Code Mapping}
\label{tab:equations}
\begin{tabular}{@{}cll@{}}
\toprule
\textbf{Eq.} & \textbf{Description} & \textbf{Code Location} \\
\midrule
1 & Router score $g(\mathbf{x})$ & \texttt{SpecificMoE.router()} \\
2 & Softmax gating $p_i$ & \texttt{F.softmax(g\_x)} \\
3 & SpecMoE Top-K output & \texttt{SpecificMoE.forward()} \\
4 & ShareMoE soft output & \texttt{SharedMoE.forward()} \\
5 & SSMoE additive fusion & \texttt{SSMoEBlock.forward()} \\
6 & L1 reconstruction loss & \texttt{pretrain\_forward()} \\
7--9 & Load-balancing loss & \texttt{load\_balancing\_loss()} \\
10 & $\mathcal{L}_\text{pretrain}$ & \texttt{pretrain\_forward()} \\
11 & $\mathcal{L}_\text{cls}$ & \texttt{FinetuneLoss} \\
\bottomrule
\end{tabular}
\end{table}

% ══════════════════════════════════════════════════════════════════════════════
% IX. DISCUSSION
% ══════════════════════════════════════════════════════════════════════════════
\section{Discussion}

\subsection{Advantages of Domain Decoupling for Landslide Mapping}

The SSMoE architecture offers a natural fit for multi-source landslide data:

\begin{enumerate}
    \item \textbf{Expert specialisation}: The Top-K routing mechanism allows different experts to specialise in different landslide-contributing factors. For instance, one expert may learn slope-failure patterns from SAR+DEM features, while another captures vegetation-loss signatures from NDVI anomalies.
    
    \item \textbf{Cross-domain interaction}: Shared experts can learn synergistic effects—e.g., the compound risk when both soil moisture \textit{and} rainfall are elevated simultaneously.
    
    \item \textbf{Computational efficiency}: Top-K routing (K=2 out of 6) activates only a subset of experts per token, reducing computational cost compared to dense FFN layers of equivalent capacity.
    
    \item \textbf{Scalability}: Additional data sources (e.g., land-use maps, geological maps, drainage networks) can be incorporated by adding channels without modifying the core architecture.
\end{enumerate}

\subsection{Heuristic Label Generation}

A key practical challenge is the absence of readily available pixel-level landslide annotations. Our heuristic label generator provides a bootstrap solution by combining DEM-derived slope (35\% weight), NDVI anomaly (30\%), soil moisture (20\%), and rainfall intensity (15\%) through a weighted composite scoring and sigmoid thresholding. While not ground-truth quality, these proxy labels enable end-to-end training and can be progressively replaced with manually digitised labels from the Landslide Atlas.

\subsection{Limitations}

\begin{itemize}
    \item \textbf{Spatial resolution mismatch}: Sentinel-2 (10~m), SMAP (9~km), and rainfall data ($0.25^{\circ}$) have vastly different native resolutions. Bilinear resampling to a common grid introduces smoothing artifacts in coarse-resolution layers.
    
    \item \textbf{Temporal alignment}: Sentinel-1/2 acquisitions may not coincide with the landslide event date. We use the closest available acquisition, which may show post-event conditions.
    
    \item \textbf{Class imbalance}: Landslide-affected pixels are a small fraction of the study area. Future work should explore focal loss or oversampling strategies.
    
    \item \textbf{Ground truth quality}: The Landslide Atlas provides approximate boundaries. Pixel-level accuracy depends on manual digitisation quality.
\end{itemize}

% ══════════════════════════════════════════════════════════════════════════════
% X. CONCLUSION
% ══════════════════════════════════════════════════════════════════════════════
\section{Conclusion}

We presented LandslideEEGMoE, a faithful adaptation of the EEGMoE domain-decoupled Mixture-of-Experts architecture for landslide susceptibility mapping. By establishing a principled mapping from EEG frequency bands to geospatial feature channels—SAR, NDVI, NDWI, rainfall, and soil moisture—we demonstrate that the SSMoE architecture's ability to learn domain-specific and cross-domain representations generalises beyond neural signal processing to multi-source remote sensing. The complete implementation includes a real geospatial data pipeline handling native satellite data formats, heuristic label generation for bootstrapping, and the exact two-stage training protocol from the original paper. This work opens avenues for applying advanced domain-decoupled architectures to other multi-source geospatial prediction tasks, including flood risk assessment, drought monitoring, and wildfire susceptibility mapping.

% ══════════════════════════════════════════════════════════════════════════════
% REFERENCES
% ══════════════════════════════════════════════════════════════════════════════
\begin{thebibliography}{00}

\bibitem{eegmoe2026}
Y.~Gao, \textit{et al.}, ``EEGMoE: A Domain-Decoupled Mixture-of-Experts Model for Self-Supervised EEG Representation Learning,'' \textit{IEEE Transactions on Neural Networks and Learning Systems}, 2026.

\bibitem{wayanad2024}
``Wayanad Landslide 2024: Impact Assessment Report,'' \textit{International Charter: Space and Major Disasters}, Charter Activation No. 1029, 2024.

\bibitem{dl_landslide}
F.~Huang, Z.~Zhang, \textit{et al.}, ``A deep learning algorithm using a fully connected sparse autoencoder neural network for landslide susceptibility prediction,'' \textit{Landslides}, vol.~17, pp.~217--229, 2020.

\bibitem{logistic_landslide}
B.~Pradhan, ``A comparative study on the predictive ability of the decision tree, support vector machine and neuro-fuzzy models in landslide susceptibility mapping,'' \textit{Computers \& Geosciences}, vol.~51, pp.~350--365, 2013.

\bibitem{svm_landslide}
H.~R.~Pourghasemi \textit{et al.}, ``Landslide susceptibility mapping using support vector machine and GIS,'' \textit{Journal of Earth System Science}, vol.~121, no.~3, pp.~681--697, 2012.

\bibitem{rf_landslide}
W.~Chen, X.~Xie, \textit{et al.}, ``A comparative study of logistic model tree, random forest, and classification and regression tree models for spatial prediction of landslide susceptibility,'' \textit{Catena}, vol.~151, pp.~147--160, 2017.

\bibitem{cnn_landslide}
Y.~Wang \textit{et al.}, ``Comparison of convolutional neural networks for landslide susceptibility mapping in Yanshan County, China,'' \textit{Science of the Total Environment}, vol.~666, pp.~975--993, 2019.

\bibitem{moe_original}
R.~A.~Jacobs, M.~I.~Jordan, S.~J.~Nowlan, and G.~E.~Hinton, ``Adaptive mixtures of local experts,'' \textit{Neural Computation}, vol.~3, no.~1, pp.~79--87, 1991.

\bibitem{switch_transformer}
W.~Fedus, B.~Zoph, and N.~Shazeer, ``Switch Transformers: Scaling to Trillion Parameter Models with Simple and Efficient Sparsity,'' \textit{Journal of Machine Learning Research}, vol.~23, pp.~1--39, 2022.

\bibitem{mae}
K.~He, X.~Chen, S.~Xie, Y.~Li, P.~Doll{\'a}r, and R.~Girshick, ``Masked Autoencoders Are Scalable Vision Learners,'' in \textit{Proc. IEEE/CVF CVPR}, 2022, pp.~16000--16009.

\end{thebibliography}

\end{document}
